%===============================================================================
% metodologia.tex
% Capítulo 3: Metodología y Arquitectura Computacional
%===============================================================================

\chapter{Metodología y Arquitectura Computacional}
\label{cap:metodologia}

La metodología seguida en este proyecto se basa en las especificaciones proporcionadas en el enunciado del curso \cite{proyecto-molina2024}.

\section{Arquitectura Híbrida: Fortran + Python}
El proyecto utiliza una arquitectura de \textit{pipeline científico} que combina dos lenguajes de programación para aprovechar sus fortalezas:

\begin{itemize}
    \item \textbf{Fortran 90}: Para el núcleo numérico y la simulación principal. Ofrece alto rendimiento en bucles matemáticos y es el estándar en computación de alto rendimiento (HPC).
    \item \textbf{Python}: Para el post-procesamiento, análisis estadístico y generación de visualizaciones. Cuenta con un ecosistema robusto de bibliotecas científicas (NumPy, Matplotlib, Pandas).
\end{itemize}

Esta separación de responsabilidades permite desacoplar la carga computacional intensiva de la fase de visualización, logrando un rendimiento óptimo.

\section{Estructura del Proyecto}
La organización del código sigue un patrón modular con responsabilidades claras:

\begin{table}[H]
    \centering
    \caption{Módulos del sistema y sus responsabilidades.}
    \label{tab:modulos}
    \begin{tabular}{ll}
        \toprule
        \textbf{Módulo} & \textbf{Responsabilidad} \\
        \midrule
        \texttt{mod\_constants.f90} & Constantes físicas y parámetros de simulación \\
        \texttt{mod\_types.f90} & Tipos de datos para el sistema de partículas \\
        \texttt{mod\_energy.f90} & Cálculo de energía electrostática y potencial \\
        \texttt{mod\_simulation.f90} & Algoritmo de minimización energética \\
        \texttt{mod\_io.f90} & Lectura de parámetros y escritura de resultados \\
        \texttt{main.f90} & Programa principal y coordinación del flujo \\
        \texttt{config.py} & Configuración de visualización y rutas \\
        \texttt{plot\_*.py} & Generación de gráficas específicas \\
        \texttt{analysis.py} & Análisis estadístico del sistema \\
        \texttt{video\_generator.py} & Generación de video de evolución \\
        \bottomrule
    \end{tabular}
\end{table}

\section{Parámetros de Simulación}
Los parámetros utilizados en las simulaciones se resumen en la tabla \ref{tab:parametros}:

\begin{table}[H]
    \centering
    \caption{Parámetros de simulación.}
    \label{tab:parametros}
    \begin{tabular}{lcc}
        \toprule
        \textbf{Parámetro} & \textbf{Simbolo} & \textbf{Valor} \\
        \midrule
        Número de partículas & \(N\) & 50 \\
        Tamaño del dominio & \(L\) & 10.0 \\
        Paso máximo de movimiento & \(\delta_{\text{max}}\) & 0.25 \\
        Iteraciones máximas & \(N_{\text{iter}}\) & 500,000 \\
        Parámetro de softening & \(\epsilon\) & 0.01 \\
        Constante de Coulomb & \(k_e\) & 1.0 (unidades reducidas) \\
        \bottomrule
    \end{tabular}
\end{table}

\section{Implementación Técnica}

\subsection{Inicialización del Sistema}
El sistema se inicializa con posiciones aleatorias uniformemente distribuidas en el dominio \([-L, L] \times [-L, L]\). Las cargas se asignan según el modo seleccionado:
\begin{itemize}
    \item \textbf{Modo 1}: Todas las cargas son \(+1\).
    \item \textbf{Modo 2}: Mezcla aleatoria de \(+1\) y \(-1\) con probabilidad 0.5 para cada carga.
\end{itemize}

\subsection{Estabilidad Numérica: Parámetro de Softening}
Para evitar la singularidad cuando dos partículas se acercan demasiado (\(r \to 0\)), implementamos un parámetro de \textit{softening} \(\epsilon\):

\begin{equation}
    r_{\text{efectivo}} = \sqrt{(x_i - x_j)^2 + (y_i - y_j)^2 + \epsilon^2}
\end{equation}

Este parámetro introduce un radio mínimo efectivo sin alterar la física a distancias \(r \gg \epsilon\).

\subsection{Formatos de Datos}
Los resultados se guardan en formato CSV (Comma-Separated Values) para garantizar compatibilidad con herramientas de análisis como Pandas y NumPy. Los archivos generados son:
\begin{itemize}
    \item \texttt{initial\_config.csv}: Configuración inicial del sistema.
    \item \texttt{final\_config.csv}: Configuración final de mínima energía.
    \item \texttt{energy\_log.csv}: Registro de energía y tasa de aceptación por iteración.
    \item \texttt{config\_XXXXXX.csv}: Configuraciones intermedias para generar el video.
\end{itemize}
